\documentclass[11pt,a4paper]{article}
\usepackage[utf8]{inputenc}
\usepackage[T1]{fontenc}
\usepackage[french]{babel}
\usepackage{lmodern}
\usepackage{amsmath,amssymb,mathtools}
\usepackage{geometry}
\usepackage{booktabs}
\usepackage{xcolor}
\usepackage{hyperref}
\geometry{margin=2.4cm}
\hypersetup{colorlinks=true,linkcolor=blue!60!black,citecolor=blue!60!black,urlcolor=blue!60!black}

\newcommand{\dd}{\,\mathrm{d}}
\newcommand{\hb}{\hbar}
\newcommand{\lp}{\ell_{\mathrm P}}
\newcommand{\tP}{t_{\mathrm P}}
\newcommand{\EP}{E_{\mathrm P}}
\newcommand{\rhoP}{\rho_{\mathrm P}}
\newcommand{\rhoL}{\rho_{\Lambda}}
\newcommand{\OmL}{\Omega_{\Lambda}}
\newcommand{\lL}{\ell_{\Lambda}}

\title{Contrainte energie noire sur un regulateur GUP de l'energie du vide}
\author{Note de verification analytique et numerique}
\date{22 mai 2026}

\begin{document}
\maketitle

\begin{abstract}
On part de la densite d'energie du vide regularisee par la mesure GUP
\[
\rho_{\mathrm{vac}}=\frac{g\rhoP}{16\pi^2\beta_0^2},
\]
et on impose formellement $\rho_{\mathrm{vac}}=\rhoL$, ou $\rhoL$ est la
densite d'energie noire deduite des parametres cosmologiques Planck 2018. Le
calcul donne une formule fermee :
\[
\beta_0^{(\Lambda)}
=
\left(\frac{g}{6\pi\OmL}\right)^{1/2}\frac{1}{H_0\tP}.
\]
Cette identification selectionne une longueur minimale effective
\[
\lL
=
\sqrt{\beta_0^{(\Lambda)}}\,\lp
=
\left(\frac{g}{6\pi\OmL}\right)^{1/4}
\sqrt{\lp\frac{c}{H_0}},
\]
c'est-a-dire le produit geometrique entre l'echelle de Planck et l'echelle de
Hubble, a un facteur numerique pres. Le resultat est mathematiquement net mais
physiquement contraignant : le meme $\beta_0$ universel rendrait les seuils
cinematiques GUP beaucoup trop bas pour les particules ordinaires.
\end{abstract}

\section{Donnees et point de depart}

On utilise les constantes CODATA/NIST 2022 et les parametres de base
$\Lambda$CDM de Planck 2018 :
\begin{equation}
H_0=67.4\ \mathrm{km\,s^{-1}\,Mpc^{-1}},
\qquad
\Omega_m=0.315,
\qquad
\OmL=1-\Omega_m=0.685.
\end{equation}

La densite d'energie noire est
\begin{equation}
\rhoL
=
\OmL\frac{3H_0^2}{8\pi G}c^2.
\label{eq:rho_lambda}
\end{equation}
Numeriquement,
\begin{equation}
\rhoL=5.253229\times10^{-10}\ \mathrm{J\,m^{-3}}.
\end{equation}

Dans le modele GUP isotrope $\gamma=3$, la contribution de point zero d'un
champ massless est
\begin{equation}
\rho_{\mathrm{vac}}
=
\frac{g\pi c}{2h^3\beta^2}.
\end{equation}
Avec
\begin{equation}
\beta=\beta_0\frac{\lp^2}{\hb^2},
\end{equation}
on obtient
\begin{equation}
\rho_{\mathrm{vac}}
=
\frac{g\rhoP}{16\pi^2\beta_0^2},
\qquad
\rhoP=\frac{c^7}{\hb G^2}.
\label{eq:rho_vac_beta0}
\end{equation}

\section{Inversion analytique}

Imposons formellement $\rho_{\mathrm{vac}}=\rhoL$. Alors
\begin{equation}
\beta_0^2
=
\frac{g\rhoP}{16\pi^2\rhoL}.
\end{equation}
En utilisant \eqref{eq:rho_lambda} et \eqref{eq:rho_vac_beta0},
\begin{equation}
\beta_0^2
=
\frac{g c^7}{16\pi^2\hb G^2}
\frac{8\pi G}{3\OmL H_0^2c^2}
=
\frac{g c^5}{6\pi\hb G\OmL H_0^2}.
\end{equation}
Comme
\begin{equation}
\tP=\sqrt{\frac{\hb G}{c^5}},
\end{equation}
on obtient la forme compacte
\begin{equation}
\boxed{
\beta_0^{(\Lambda)}
=
\left(\frac{g}{6\pi\OmL}\right)^{1/2}
\frac{1}{H_0\tP}.
}
\label{eq:beta0_hubble_planck}
\end{equation}

Pour $g=1$,
\begin{equation}
\beta_0^{(\Lambda)}=2.363228\times10^{60}.
\end{equation}
Ainsi, le resultat numerique $10^{60}$ n'est pas accidentel : il vient
directement de l'enorme rapport entre le temps de Hubble et le temps de Planck.

\section{Echelle de longueur emergente}

La longueur minimale effective correspondant a
\eqref{eq:beta0_hubble_planck} est
\begin{equation}
\lL=\sqrt{\beta_0^{(\Lambda)}}\,\lp.
\end{equation}
Donc
\begin{equation}
\boxed{
\lL
=
\left(\frac{g}{6\pi\OmL}\right)^{1/4}
\sqrt{\lp\frac{c}{H_0}}.
}
\label{eq:ell_lambda}
\end{equation}

Cette formule est l'une des consequences les plus utiles du test. Elle montre
que l'identification directe de $\rho_{\mathrm{vac}}$ a $\rhoL$ ne selectionne
pas l'echelle de Planck, mais une echelle UV/IR mixte :
\begin{equation}
\lL\propto\sqrt{\lp R_H},
\qquad
R_H=\frac{c}{H_0}.
\end{equation}

Pour $g=1$,
\begin{equation}
\lL=2.484636\times10^{-5}\ \mathrm{m}.
\end{equation}

\section{Echelle d'energie associee}

L'energie GUP associee a cette longueur est
\begin{equation}
E_{\mathrm{GUP}}=\frac{\hb c}{\lL}.
\end{equation}
En utilisant \eqref{eq:ell_lambda},
\begin{equation}
\boxed{
E_{\mathrm{GUP}}^{(\Lambda)}
=
\left(\frac{6\pi\OmL}{g}\right)^{1/4}
\sqrt{\EP\,\hb H_0}.
}
\label{eq:energy_geomean}
\end{equation}

L'energie obtenue est donc le produit geometrique entre l'energie de Planck
$\EP$ et l'energie de Hubble $\hb H_0$, a un facteur numerique pres. Pour
$g=1$,
\begin{equation}
E_{\mathrm{GUP}}^{(\Lambda)}=7.94\ \mathrm{meV}.
\end{equation}

\section{Contrainte cinematique croisee}

La reformulation position--vitesse donne
\begin{equation}
v_{\mathrm{GUP}}
=
c\frac{m_P}{m\sqrt{\beta_0}}.
\end{equation}
Sous l'identification a l'energie noire,
\begin{equation}
\boxed{
\frac{v_{\mathrm{GUP}}}{c}
=
\frac{m_*}{m},
\qquad
m_*c^2=E_{\mathrm{GUP}}^{(\Lambda)}.
}
\label{eq:v_threshold_mass_star}
\end{equation}

Pour $g=1$,
\begin{equation}
m_*c^2=7.94\ \mathrm{meV}.
\end{equation}
Ainsi, une particule de masse superieure a cette echelle aurait un seuil GUP
subluminal si le meme $\beta_0$ etait applique universellement.

\begin{table}[h]
\centering
\begin{tabular}{lcc}
\toprule
Particule & $v_{\mathrm{GUP}}/c$ avec $\beta_0^{(\Lambda)}$ & Interpretation \\
\midrule
electron & $1.55\times10^{-8}$ & exclut une deformation universelle simple \\
proton & $8.46\times10^{-12}$ & exclut une deformation universelle simple \\
neutrino $0.05$ eV & $1.59\times10^{-1}$ & deja subluminal \\
\bottomrule
\end{tabular}
\caption{Seuils vitesse obtenus si l'on applique a la cinematique des particules le $\beta_0$ ajuste a l'energie noire pour $g=1$.}
\end{table}

\section{Densite maximale de modes ajustee a $\rhoL$}

Le premier manuscrit donne
\begin{equation}
n_{\max}=\frac{g}{32\pi\beta_0^{3/2}\lp^3}.
\end{equation}
En inserant \eqref{eq:beta0_hubble_planck},
\begin{equation}
\boxed{
n_{\max}^{(\Lambda)}
=
\frac{g^{1/4}(6\pi\OmL)^{3/4}}{32\pi}
\frac{1}{(\lp R_H)^{3/2}}.
}
\label{eq:nmax_lambda}
\end{equation}
La distance moyenne entre modes vaut donc
\begin{equation}
\boxed{
\left(n_{\max}^{(\Lambda)}\right)^{-1/3}
=
\left[\frac{32\pi}{g^{1/4}(6\pi\OmL)^{3/4}}\right]^{1/3}
\sqrt{\lp R_H}.
}
\label{eq:spacing_lambda}
\end{equation}

Pour $g=1$,
\begin{equation}
n_{\max}^{(\Lambda)}=6.49\times10^{11}\ \mathrm{m^{-3}},
\qquad
\left(n_{\max}^{(\Lambda)}\right)^{-1/3}=1.16\times10^{-4}\ \mathrm{m}.
\end{equation}

\section{Conclusion}

Le test donne une conclusion mathematique claire. Une mesure GUP peut
regulariser l'energie du vide et produire une valeur finie. Mais si l'on impose
que cette valeur soit exactement l'energie noire observee, alors
\begin{equation}
\beta_0\sim\frac{1}{H_0\tP}\sim10^{60}.
\end{equation}
Cette valeur introduit une longueur minimale effective microscopique mais non
planckienne, de l'ordre de $10^{-5}$ m, et une energie GUP de l'ordre du meV.

La consequence physique est restrictive : une telle valeur ne peut pas etre le
meme parametre universel qui controle la cinematique des particules ordinaires.
Le modele doit donc choisir l'une des options suivantes :
\begin{enumerate}
\item traiter la regularisation du vide comme un secteur effectif distinct ;
\item introduire une annulation ou une renormalisation gravitationnelle ;
\item rendre $\beta_0$ dependant du secteur ou de l'etat ;
\item modifier la mesure par une construction covariante plus riche.
\end{enumerate}

\begin{thebibliography}{9}
\bibitem{CODATA2022}
E. Tiesinga, P. J. Mohr, D. B. Newell and B. N. Taylor,
``CODATA recommended values of the fundamental physical constants: 2022'',
\emph{Reviews of Modern Physics} \textbf{97}, 025002 (2025).

\bibitem{Planck2018}
N. Aghanim et al. (Planck Collaboration),
``Planck 2018 results. VI. Cosmological parameters'',
\emph{Astronomy \& Astrophysics} \textbf{641}, A6 (2020),
arXiv:1807.06209.
\end{thebibliography}

\end{document}
